\section{Introduction}
Image enhancement is the process of improving an image's quality and information content~\cite{singh2014various}. The goal is to increase visual differences among image features, making it more suitable for various applications, such as increasing the brightness of dark images for better viewing. Common image enhancement techniques include sharpening, smoothing, contrast enhancement, and noise reduction~\cite{draa2014artificial}. Contrast enhancement, in particular, is applied to images or videos to increase their dynamic range~\cite{hashemi2010image}.

Image enhancement has been practiced as an applicable problem in meta-heuristics for a long time~\cite{cuckoo_search_and_PSO,ACO_GA_SA,munteanu2004gray,mukhopadhyay2023image}. Moreover, the integration of fuzzy systems with meta-heuristic algorithms has been the subject of area within the Computational Intelligence 
community~\cite{fuzzy_plus_meta,diaz2017new}. Many studies have approached image enhancement as an optimization problem, focusing on modifying the image quality fitness function~\cite{cuckoo_search_and_PSO,draa2014artificial,oloyede2019new}, combining several meta-heuristic algorithms~\cite{cuckoo_search_and_PSO,ACO_GA_SA}, optimizing parameters \cite{samanta2018log} and avoiding local optima \cite{dos2009differential}. However, one aspect that has not received as much attention in image enhancement problems is the adaptation of fuzzy systems. Sandeep and Samrudh~\cite{s2018image} have demonstrated that the variation in the input membership function in a fuzzy system can positively impact image enhancement performance.

This paper addresses the image contrast enhancement problem by using stochastic optimization within a fuzzy logic system. 
Our primary contribution is the design of a meta-heuristic inspired framework that leverages a fuzzy logic system.
The core component of the fuzzy logic system is a set of {\em input membership function} that are used to describe an {\em intensity transformation function}. 
We implement genetic operators that {\em mutate} the fuzzy logic system, thereby enhancing the original image to achieve optimal contrast enhancement. 
The fuzzy system in our approach is based on simple contrast enhancement rules, outlined in \cite{gonzalez2008digital}, and it operates independently of the input image.

Our main idea of the fuzzy image enhancement technique is as follows: we start with a basic fuzzy rule set and a set of input membership functions. By applying a metaheuristics framework, we try to evolve the fuzzy sets. Finally, we apply the transformation function described by fuzzy sets on the value channel of HSV color space to get the final image. Thus we have converted the problem into an optimization problem. In other words, rather than generating an image, we try to generate a suitable mapping between the input and output color values. However, this is inherently challenging since, even if we consider only 8-bit color, i.e., $256$ color values, the number of possible mappings becomes huge. Thus it becomes infeasible to search through the solution space exhaustively, and there is no definite knowledge about how to improve/generate a solution. This motivates us to leverage a metaheuristics framework \cite{Luke2009Metaheuristics}.